% 《大学物理简明教程》（第五版）第1章习题：
% 1.1.1--1.1.7（印刷/PDF第15--16页）
% 1.2.1--1.2.7（印刷/PDF第16--17页）
% 1.3.1--1.3.3（印刷/PDF第17页）
% 按扫描本独立推导；题干为保留条件的研究用转述。质量状态：待独立复核。

\subsection{选择题：1.1.1--1.1.7}
\sourceproblem{1.1.1}{15}{15}
\problemtext 已知质点的位置矢量$\bm r(t)=(x(t),y(t))$，判断瞬时速度大小的正确表达式。
\workings
定义速度矢量为$\bm v=d\bm r/dt=(\dot x,\dot y)$，故
\[
 \boxed{v=|\bm v|=\sqrt{\dot x^{\,2}+\dot y^{\,2}}}.
\]
应选原题\textbf{D}。$\dot{\bm r}$是矢量，而$\frac{d|\bm r|}{dt}$是径向速度，可与实际速率不同。
\crosscheck 匀速圆周运动中$|\bm r|$恒定，其导数为零但$v>0$，排除以径向变化率代替速率的选项。

\sourceproblem{1.1.2}{16}{16}
\problemtext 一质点沿直线运动，某时刻\emph{速率}为$2\,\mathrm{m/s}$，有符号加速度$a=-2\,\mathrm{m/s^2}$，询问$1$秒后的速率。
\workings
若初始速度为$+2\,\mathrm{m/s}$，则$v(1)=2-2=0$；若初始速度为$-2\,\mathrm{m/s}$，则$v(1)=-2-2=-4\,\mathrm{m/s}$，速率为$4\,\mathrm{m/s}$。
仅给\emph{速率}，未给初始运动方向，不能确定唯一结果；原题选\textbf{D}。
\crosscheck 速率非负，不能将带符号加速度直接从速率中代数相减而忽略初始速度符号。

\sourceproblem{1.1.3}{16}{16}
\problemtext 质点在半径$R$的圆周上做匀速率运动，每$t$秒转一整圈；在$2t$内求平均速度的大小和平均速率。
\workings
$2t$内恰走两圈，始末点重合，故平均速度大小为
$|\bar{\bm v}|=|\Delta\bm r|/(2t)=0$。
总路程为$4\pi R$，平均速率为$4\pi R/(2t)=2\pi R/t$。
\[
\boxed{|\bar{\bm v}|=0,\qquad \bar v_{\mathrm{speed}}=\frac{2\pi R}{t}}.
\]
原题选\textbf{B}。
\crosscheck 平均速度是矢量，平均速率是标量，不能混淆。

\sourceproblem{1.1.4}{16}{16}
\problemtext 对曲线运动的位矢$\bm r$、速度$\bm v$、加速度$\bm a$、路程$s$和切向加速度$a_\tau$，考查以下导数关系是否成立。
\workings
严格的矢量、标量恒等式是
\[
 \boxed{\frac{d\bm r}{dt}=\bm v,\quad
 \frac{d\bm v}{dt}=\bm a,\quad
 \frac{ds}{dt}=|\bm v|,\quad
 a_\tau=\frac{d|\bm v|}{dt}.}
\]
因此$\left|\frac{d\bm v}{dt}\right|=|\bm a|$，一般不等于$|a_\tau|$；
只有法向加速度为零的情形才可能相等。
\crosscheck 原书四种写法的标量$v$、矢量$\bm v$在扫描字体中较接近，
若不辨明粗体，就无法严谨判定题面的①、③、④及给定选项；本题\textbf{暂不强行判定唯一选项}，待逐字二校。

\sourceproblem{1.1.5}{16}{16}
\problemtext 平面一般曲线运动中，以$v$表示瞬时速率、$\bar v$表示某时间区间的平均速率，以$\bm v,\bar{\bm v}$表示相应速度矢量，判断它们的模是否等于标量速率。
\workings
按定义$|\bm v(t)|=v(t)$始终成立。另一方面，
\[
 |\bar{\bm v}|=\frac{|\Delta\bm r|}{\Delta t}
 \le \frac{\text{路程}}{\Delta t}=\bar v,
\]
一般曲线运动且在区间内发生转向时严格小于，故原题选\textbf{D}。
\crosscheck 当轨迹退化为方向始终不变的直线时，两种平均值可相等；“一般曲线”不是绝对严格不等号的逻辑充分条件，应理解为通常不等。

\sourceproblem{1.1.6}{16}{16}
\problemtext 一维位置$x=3t-5t^3+6$（SI），选择运动类型及加速度方向。
\workings
\[
 v=3-15t^2,\qquad a=\frac{dv}{dt}=-30t.
\]
在$t>0$，加速度沿$x$轴负方向，其数值随时间变化，属于变加速直线运动，原题选\textbf{D}。
\crosscheck 加速度朝负方向\emph{不}保证速率持续减小；$t>1/\sqrt5$时速度也为负，速率反而增加。$t=0$时加速度恰为零。

\sourceproblem{1.1.7}{16}{16}
\problemtext 判断以下命题：沿$x$轴运动若$a<0$就必减速；曲线运动的加速度每一时刻均不为零；加速度恒矢量就必为直线运动；抛体运动切/法向加速度变化则合加速度大小必变化。
\workings
逐项构造反例：
负向速度$v<0$且$a<0$时速率\emph{增加}；
非直线轨迹可以在孤立时刻有$a=0$（例如$\bm r(t)=(t,t^3)$在$t=0$）；
斜抛运动有恒定$\bm a=-g\bm j$却不是直线；
理想抛体的切向与法向分量随方向变化，但合加速度大小恒为$g$。
因此按题面四个命题\textbf{均不成立}。
\crosscheck 此题的原版选项未列“以上均不正确”，疑有命题或印刷问题，保留\textbf{待核实}状态，不凭空指定字母答案。

\subsection{填空题：1.2.1--1.2.7}
\sourceproblem{1.2.1}{16}{16}
\problemtext 半径$R=5\,\mathrm m$，速率$v=\pi\,\mathrm{m/s}$匀速圆周运动，求$5\,\mathrm s$内位移大小和路程。
\workings
角位移$\Delta\theta=vt/R=\pi$，运动半周，
\[
 \boxed{|\Delta\bm r|=2R=10\,\mathrm m,\qquad
 s=vt=5\pi\,\mathrm m}.
\]
\crosscheck 半周位移为直径，不是弧长。

\sourceproblem{1.2.2}{16}{16}
\problemtext 一维加速度$a(t)=3+2t$（SI），初速率$v_0=5\,\mathrm{m/s}$，求$t=3\,\mathrm s$速率（初速度取与加速度同方向）。
\workings
$v(t)=v_0+\int_0^t(3+2u)\,du=5+3t+t^2$，故
\[
\boxed{v(3)=23\,\mathrm{m/s}}.
\]
\crosscheck $a(t)>0$且$v_0>0$，全程不转向，速度和速率一致。

\sourceproblem{1.2.3}{16}{16}
\problemtext 质点从静止开始做$R=1\,\mathrm m$圆周运动，角加速度$\alpha=12t^2-6t$（SI），求$\omega(t)$及切向加速度大小。
\workings
\[
 \omega(t)=\int_0^t(12u^2-6u)\,du
 =4t^3-3t^2,\qquad
 \boxed{a_\tau=R|\alpha|=|12t^2-6t|\,\mathrm{m/s^2}}.
\]
\crosscheck 切向加速度的\emph{大小}不可写成可能为负的$\alpha$；$\omega(0)=0$。

\sourceproblem{1.2.4}{16--17}{16--17}
\problemtext 图1.2.4给出的$x$--$t$关系是向下凹的拱形曲线：$t=0,6\,\mathrm s$时$x=0$，在$t=3\,\mathrm s$达到最大值$5\,\mathrm m$。
求瞬时速度为零时刻和速度、加速度同向的时间区间。
\workings
位置曲线在峰顶切线水平，所以$\boxed{t=3\,\mathrm s}$时$v=\dot x=0$。
图中斜率在$0<t<3$时正，在$3<t<6$时负；
曲线始终下凹，$\ddot x<0$。
因此在$\boxed{3<t<6\,\mathrm s}$，$v<0,a<0$，两者同向。
\crosscheck 由图读到的是速度斜率及加速度凹凸性，不需要擅自把原图拟合成唯一抛物线。

\sourceproblem{1.2.5}{17}{17}
\problemtext 一维质点的速率与路程关系$v=1+s^2$（SI），求任意$s$处切向加速度大小。
\workings
链式法则$\frac{dv}{dt}=\frac{dv}{ds}\frac{ds}{dt}=v\,dv/ds$，所以
\[
\boxed{a_\tau=2s(1+s^2)}\quad\text{(SI数值形式)}.
\]
\crosscheck $s\ge0$且$v>0$，故式值非负；$s=0$时$a_\tau=0$。

\sourceproblem{1.2.6}{17}{17}
\problemtext 平面运动$\bm r=2t\,\bm i+(2-t^2)\,\bm j$（SI）。求$1$至$2\,\mathrm s$的位移、$t=2\,\mathrm s$的加速度，以及轨迹方程。
\workings
\[
\Delta\bm r=\bm r(2)-\bm r(1)
 =(4\bm i-2\bm j)-(2\bm i+\bm j)
 =\boxed{2\bm i-3\bm j\ \mathrm m}.
\]
\[
\bm v=2\bm i-2t\bm j,\quad
\boxed{\bm a(2)=-2\bm j\ \mathrm{m/s^2}}.
\]
消去$t=x/2$得$\boxed{y=2-x^2/4}$（$x,y$以米的数值表示）。
\crosscheck $\bm a$与$t$无关，轨迹为向下开口抛物线。

\sourceproblem{1.2.7}{17}{17}
\problemtext 跳水运动员从$10\,\mathrm m$高处自由落下，入水后受阻减速，沿入水向下方向满足$a=-kv^2$，$k=0.4\,\mathrm{m^{-1}}$；入水速度减为初值的$10\%$时，求入水深度。
\workings
取入水点为$s=0$，以向下为正。水中$a=v\,dv/ds$，
\[
v\frac{dv}{ds}=-kv^2 \ \Longrightarrow\
\frac{dv}{v}=-k\,ds,\qquad
v(s)=v_{\mathrm{in}}e^{-ks}.
\]
当$v(s)=0.1v_{\mathrm{in}}$时
\[
\boxed{s=\frac{\ln10}{k}=2.5\ln10\ \mathrm m\approx5.756\,\mathrm m}.
\]
\crosscheck 初始落差$10\,\mathrm m$只决定$v_{\mathrm{in}}=\sqrt{20g}$；在纯$v^2$阻力模型中速度比与$v_{\mathrm{in}}$约去。此式将题目给出的水中\emph{总加速度}视为$-kv^2$，未另外加重力。

\subsection{解答题：1.3.1--1.3.3}
\sourceproblem{1.3.1}{17}{17}
\problemtext 在四条一维位置曲线$x_1=4t-3$、$x_2=-4t^3+3t^2+6$、$x_3=-2t^2+8t+4$及$x_4=2/t^2-4/t$中找出匀变速直线运动；求其$t=3\,\mathrm s$的速度、加速度并判断增速或减速。
\workings
匀变速要求$\ddot x$为常量。逐项求导：
\[
\begin{array}{l|l|l}
x(t)&v=\dot x&a=\ddot x\\ \hline
x_1&4&0\\
x_2&-12t^2+6t&-24t+6\\
x_3&-4t+8&-4\\
x_4&-4t^{-3}+4t^{-2}&12t^{-4}-8t^{-3}
\end{array}
\]
所以若“匀变速”包含$a=0$的退化情形，(1)、(3)属于恒加速度；
若特别指\emph{非零}恒加速度，则仅(3)。
$t=3$时(1)有$(v,a)=(4,0)$，匀速；
(3)有$(v,a)=(-4,-4)$（SI），同向故\emph{速率增加}。
\crosscheck (4)在$t=0$没有定义，不能称为全实数时间上的规则运动。

\sourceproblem{1.3.2}{17}{17}
\problemtext 对匀速直线、匀速曲线、变速直线和变速曲线四种运动，分别判断质点的\emph{速度矢量}和\emph{加速度矢量}是否为零。
\workings
匀速直线：$v\ne0$，$a=0$。
匀速曲线：$v\ne0$，若曲率$\kappa>0$则$a_n=v^2\kappa>0$。
变速直线：$v$可能在转向瞬间为零，非零$\dot v$时$a\ne0$。
变速曲线：$\bm v$可在孤立时刻为零，$\bm a$可在孤立时刻为零；不能仅凭“曲线”或“变速”字样断言\emph{每一瞬时}非零。
\crosscheck 非零速率的平滑弯曲段满足$a_n>0$，但在拐点可以有曲率零；静止瞬间需分析具体轨迹的可微性。

\sourceproblem{1.3.3}{17}{17}
\problemtext 给定平面坐标$x(t),y(t)$，比较两种求速度、加速度的方法：
先求径向距离$r=\sqrt{x^2+y^2}$再求导，与先分别求坐标分量再合成矢量模。何者正确，为什么？
\workings
\[
\bm v=\dot x\,\bm i+\dot y\,\bm j,\quad
v=\sqrt{\dot x^2+\dot y^2},\quad
\bm a=\ddot x\,\bm i+\ddot y\,\bm j,\quad
a=\sqrt{\ddot x^2+\ddot y^2}.
\]
而
\[
\dot r=\frac{x\dot x+y\dot y}{\sqrt{x^2+y^2}}
=\bm v\cdot\hat{\bm r}
\]
只是速度的\emph{径向分量}，且
$\ddot r$还包含曲率相关项，通常不等于加速度的模。
正确方法为先微分\emph{位置矢量}的各分量，再取模；
将\emph{位置矢量}微分本身当然也正确，但不是先微分其模。
\crosscheck 匀速圆周运动$r=R$使$\dot r=\ddot r=0$，而$v=R\omega$和$a=R\omega^2$均可非零。
